\bookchapter{الفصل 68: طريق التقدم — الأذونات والإشعارات والمعماريات الحديثة}

\chapterintro{صناعة التطبيق بداية فقط. نتحدث في هذا الفصل عن الحواجز بين «أن تعرف كيف تصنع» و«أن تصنع جيدًا»: كيف تطلب الأذونات بلطف، وكيف ترسل إشعارات لا تثير السخط، وكيف ترفع المعمارية درجة أخرى، وكيف — في عام 2026 — تشرك الذكاء الاصطناعي في كتابة كود iOS. وكل بند منها ستصادفه قريبًا في مشاريع حقيقية.}

\booksection{أولًا: الأذونات — اطرق الباب أولًا ثم ادخل}

الكاميرا والموقع وجهات الاتصال والإشعارات… هذه «المناطق الحساسة» يطلب iOS فيها من التطبيق أن يستأذن المستخدم بنافذة منبثقة عند الحاجة إليها. والقاعدة ثلاث خطوات: «التصريح + الطلب + التدهور الرشيق»:

\begin{lstlisting}[language=Swift]
// 1. صرّح في Info.plist بنص الغرض (بدون تصريح سينهار التطبيق فورًا!)
// Privacy - Location When In Use Usage Description: "يُستخدم للاستعلام عن طقس مدينتك"

// 2. طلبه عند الاستخدام
import CoreLocation

@Observable
class LocationModel {
    private let manager = CLLocationManager()

    func request() {
        switch manager.authorizationStatus {
        case .notDetermined: manager.requestWhenInUseAuthorization()
        case .denied, .restricted: /* التوجيه إلى الإعدادات لتفعيله */
        default: /* الإذن ممنوح، انطلق */
        }
    }
}

// 3. تدهور رشيق بعد الرفض: عرض تنبيه والانتقال إلى الإعدادات
// UIApplication.shared.open(URL(string: UIApplication.openSettingsURLString)!)
\end{lstlisting}

\begin{tcbwarn}{تنبيه: أذونات iOS «تطلب مرة واحدة»}
نافذة إذن iOS تظهر مرة واحدة فقط — فبعد رفض المستخدم لن تظهر أبدًا. لذلك: اختر لحظة الطلب عند «اللحظة التي يحتاج فيها المستخدم الميزة فعلًا» (لا عند الإقلاع مباشرة بمجموعة نوافذ)، وليكن النص يوضح الغرض. وتوجيه الواجهة بعد الرفض («افتح الموقع من الإعدادات») أمر أساسي، فلا تحبس المستخدم في طريق مسدود.

\end{tcbwarn}

\booksection{ثانيًا: الإشعارات — قل كلامًا نافعًا في الوقت الصحيح}

الإشعارات قناة وصول التطبيق إلى المستخدم، وهي أيضًا أكثر الوظائف عرضة لأن يضعها المستخدم في «القائمة السوداء». يستخدم iOS إطار UserNotifications، ولإرسال الإشعارات ثلاث قواعد: أرسل المفيد فقط، واسمح للمستخدم بإيقافها، ولا تزعج أحدًا في منتصف الليل. والاستخدام الأساسي:

\begin{lstlisting}[language=Swift]
import UserNotifications

// 1. طلب إذن الإشعارات (نافذة منبثقة مرة واحدة)
UNUserNotificationCenter.current()
    .requestAuthorization(options: [.alert, .sound, .badge]) { granted, _ in
        print(granted ? "ممنوح" : "مرفوض")
    }

// 2. جدولة "تذكير يومي الساعة 20:00 لتسجيل المصروفات"
func scheduleDailyReminder() {
    let content = UNMutableNotificationContent()
    content.title = "تذكير بتسجيل المصروفات"
    content.body = "لم تسجل أي مصروفات اليوم بعد"
    content.sound = .default

    var comps = DateComponents()
    comps.hour = 20
    let trigger = UNCalendarNotificationTrigger(dateMatching: comps, repeats: true)

    let request = UNNotificationRequest(
        identifier: "daily-reminder",
        content: content,
        trigger: trigger
    )
    UNUserNotificationCenter.current().add(request)
}
\end{lstlisting}

انتبه: إذن الإشعارات في iOS كذلك «يطلب مرة واحدة»، ويمكن للمستخدم إيقافه في أي وقت من إعدادات النظام — لذا فالإحساس بقيمة محتوى الإشعار مهم جدًا؛ فإرسال عشوائي واحد قد يفقدك المستخدم إلى الأبد.

\booksection{ثالثًا: ترقية المعمارية — تعميق MVVM وقابلية الاختبار}

استخدمنا MVVM في الفصل 65، ونضيف هنا ثلاث نقاط متقدمة:

\begin{itemize}
  \item التجريد بالبروتوكول: اجعل ViewModel يعتمد على بروتوكول (protocol) لا على صنف محدد — فتستطيع في الاختبار حقن «مصدر بيانات وهمي»، وهذا حجر الأساس في قابلية الاختبار؛
  \item المسؤولية الواحدة: ViewModel واحد يتكفل بمهمة واحدة (المسؤول عن الحسابات لا يتولى الطقس)؛ وإذا كثرت الوظائف ففككه؛
  \item لا تبالغ في التصميم: في مشروع صغير لشخص واحد يكفي «صنف @Observable مع تجريد بالبروتوكول». أما المعماريات الثقيلة مثل VIPER وClean Architecture فانتظر حتى يطرق بابك فريق عمل وتعقيد حقيقي — فالأدوات تخدم الأهداف، ولا تلبس للتفاخر.
\end{itemize}

\begin{tcbinfo}{نصيحة: «مكافآت البيئة» في SwiftUI}
في بيئة iOS أيضًا أسلوبان متقدمان «لا بديل عنهما»: الأدوات المصغرة (WidgetKit) تتيح للمستخدم رؤية المعلومات دون فتح التطبيق (كوضع «مصروفات هذا الشهر» من الدفتر على سطح المكتب)، والأوامر السريعة (App Intents) تتيح للمستخدم استدعاء وظائفك عبر Siri أو الأتمتة. وكلاهما له إطار رسمي وتكلفة تعلم ليست عالية، لكنهما يرفعان تجربة المستخدم ارتفاعًا هائلًا — فالأولوية على طريق التقدم أن تفتحهما أولًا.

\end{tcbinfo}

\booksection{رابعًا: أسلوب 2026 الجديد — تطوير iOS بمساعدة الذكاء الاصطناعي}

في عام 2026 صار الذكاء الاصطناعي مدمجًا بعمق في دورة تطوير Apple. وإن أردت التسريع فهذه التوليفة ناضجة جدًا:

\begin{center}
\begin{tabularx}{\linewidth}{|X|X|X|}
\hline
\textbf{المرحلة} & \textbf{كيف يساعد الذكاء الاصطناعي} & \textbf{مسؤوليتك} \tabularnewline
\hline
كتابة الكود & إكمال الكود المدمج في Xcode + توليد عبر إضافات ذكاء اصطناعي من طرف ثالث & وصف المتطلبات بوضوح ومراجعة الناتج \tabularnewline
البحث عن المعلومات & إعطاء الوثائق الرسمية ورسائل الخطأ إلى الذكاء الاصطناعي ليشرحها & التحقق من النتائج المهمة بالعودة إلى الوثائق الرسمية \tabularnewline
كتابة الاختبارات & اطلب من الذكاء الاصطناعي توليد حالات XCTest بثلاث فئات: «عادي / حدّي / استثنائي» & التأكد من أن الحالات تنسجم مع دلالات مجال العمل \tabularnewline
إصلاح الأخطاء & لصق سجل الانهيار مع الكود ذي الصلة، والذكاء الاصطناعي يحدد الموضع ويصلح & فهم التعديلات وتشغيل الاختبارات للتحقق \tabularnewline
\hline
\end{tabularx}
\end{center}

\begin{center}\small\color{gray}المشاهد الأربعة لتطوير iOS بمساعدة الذكاء الاصطناعي\end{center}

أما المنهجية التفصيلية — كيف توجه الذكاء الاصطناعي بالتعليمات، وكيف تقبل العمل، وكيف تتجنب الفخاخ — فهناك 11 فصلًا كاملة عنها في الجزء السابع من هذا الكتاب «البرمجة العملية بالذكاء الاصطناعي» (الفصول 37-47)، وننصح بشدة بقراءتها معًا: اترك قواعد Swift لهذا الجزء، وأسرار التعاون مع الذكاء الاصطناعي للجزء السابع.

\begin{tcbpractice}{هيا نجرّب: أضف «الثالوث» إلى دفتر الحسابات}
\booksection{خامسًا: أساسيات التزامن في Swift — @MainActor وTask}

في منظومة التزامن في Swift 6 قاعدتان يجب على كل مبتدئ فهمهما:

\begin{itemize}
  \item الواجهة لا تُحدَّث إلا على الخيط الرئيسي: عروض SwiftUI من نوع @MainActor افتراضيًا، فعند عودة نتيجة الشبكة ارجع إلى الخيط الرئيسي ثم غيّر الحالة (عبر await MainActor.run \{ ... \} أو بوسم ViewModel بـ @MainActor)؛
  \item المهام غير المتزامنة تديرها Task: المُعدِّل .task \{ \} يبدأ تلقائيًا مع ظهور العرض ويُلغى تلقائيًا عند اختفائه — وهذه هي الطريقة المعيارية لبدء عمل غير متزامن في SwiftUI، فلا تكتب Task عشوائيًا داخل body مباشرة.
\end{itemize}

\begin{lstlisting}[language=Swift]
struct WeatherScreen: View {
    @State private var vm = WeatherViewModel()

    var body: some View {
        Group {
            switch vm.state {
            case .loading: ProgressView()
            case .success(let d): WeatherContent(data: d)
            case .error(let m): Text(m)
            }
        }
        .task {                       // يبدأ عند الظهور ويُلغى عند الاختفاء
            await vm.search(city: "الرياض")
        }
    }
}
\end{lstlisting}

خلف هاتين القاعدتين تصميم «أمان سباقات البيانات» في Swift 6 — فالمترجم يخبرك بالخطأ عند كتابته مباشرة، بدلاً من اكتشافه بعد انهيار التطبيق. وخطأ المترجم ليس تنمرًا، بل هو يتقاضى عنك الرصاصات.

(1) أضف إلى الدفتر «تذكيرًا يوميًا الساعة 20:00 لتسجيل الحساب» (إذن الإشعارات + إشعار محلي)؛ (2) حوّل مصدر بيانات ViewModel إلى تجريد بالبروتوكول، واكتب اختبار وحدات يحقن بيانات وهمية؛ (3) أضف أداة مصغرة على سطح المكتب لـ«مصروفات هذا الشهر». وبإتمام الثلاثة يكتسب مشروعك هيئة «الجيش النظامي».

\end{tcbpractice}

\begin{tcbsummary}{خلاصة الفصل}
\begin{itemize}
  \item الأذونات في ثلاث خطوات: التصريح في Info.plist (فإذا أُهمل انهار التطبيق)، والطلب عند الاستخدام، وتوجيه المستخدم إلى الإعدادات عند الرفض
  \item نافذة إذن iOS تظهر مرة واحدة: اختر لحظة «الاستخدام الفعلي»، وليكن النص واضح الغرض
  \item قواعد الإشعارات الثلاث: المفيد فقط، ويمكن إيقافها، ولا إزعاج؛ عبر إطار UserNotifications
  \item ترقية MVVM: تجريد بالبروتوكول (قابل للاختبار)، ومسؤولية واحدة، وبلا مبالغة في التصميم
  \item مكافآت البيئة: أدوات مصغرة بـ WidgetKit + أوامر سريعة بـ App Intents
  \item أربعة مشاهد لمساعدة الذكاء الاصطناعي: كتابة الكود، والبحث عن المعلومات، وكتابة الاختبارات، وإصلاح الأخطاء؛ والمنهجية في الجزء السابع (الفصول 37-47)
\end{itemize}

\end{tcbsummary}
